\section{Conclusion}
\label{sec:conclusion}

This research provides a strictly controlled empirical evaluation of three major exact solving paradigms---Boolean SAT, Constraint Programming, and Mixed Integer Programming---on the N-Queens problem. Through the execution of 495 scheduled benchmark trials (430 successful completions) up to $N=100$, we observed that theoretical SAT structural compactness is a poor predictor of practical solving efficiency. Massive clause models (Pairwise) frequently outpaced compact-clause implementations (Sequential Counter). Under default heuristic settings, the SAT-Binary encoding emerged as the most efficient boolean model, completing $N=100$ in 0.231 seconds.

At $N=200$, the limitations of both hardware and software environments became evident. The SAT-Pairwise encoding was excluded due to an estimated memory footprint exceeding hardware limits, while IBM CP Optimizer encountered a license-imposed search space restriction. Among the successful methods at $N=200$, SAT-Commander achieved the fastest median resolution time of 0.441 seconds, significantly outpacing OR-Tools CP-SAT (3.672 seconds) and suggesting that optimal boolean partitioning can remain highly effective even at massive scales.
However, native Constraint Programming implementations handled the large state spaces efficiently. The IBM CP Optimizer recorded the fastest runtime overall, resolving $N=100$ in 0.153 seconds by circumventing the massive memory initialization bottlenecks required by SAT translation. Nevertheless, CP frameworks still suffer from exponential backtracks and do not universally eliminate search limitations. Furthermore, we documented severe non-monotonic search pathologies in the SAT-Product encoding, underscoring the extreme sensitivity of CDCL solvers to variable layout. 

Future investigations should focus on phase-policy ablation studies to isolate the mechanisms triggering SAT search pathologies, as well as evaluating these structural encodings on a broader suite of combinatorial benchmarks beyond placement constraints.
